%%=============================================================================
%% Voorwoord
%%=============================================================================

\chapter*{\IfLanguageName{dutch}{Woord vooraf}{Preface}}%
\label{ch:voorwoord}

%% TODO:
%% Het voorwoord is het enige deel van de bachelorproef waar je vanuit je
%% eigen standpunt (``ik-vorm'') mag schrijven. Je kan hier bv. motiveren
%% waarom jij het onderwerp wil bespreken.
%% Vergeet ook niet te bedanken wie je geholpen/gesteund/... heeft
This bachelor thesis focuses on automatic technical feedback on the kip on bars through video analysis and Large Language Models. I chose this topic because I personally enjoy sports and find it interesting how technology can support athletes and coaches. Gymnastics is also a sport that I find technically very impressive. The combination of movement, timing, strength and control made the kip on bars an interesting skill for me to investigate further.

The process was not always straightforward. Some ideas turned out to be technically more difficult than expected, and certain results were less strong than hoped. However, these limitations helped to make the proof-of-concept more realistic and better defined.

I would like to thank my supervisor, Ms. K. Samyn, for the regular feedback and guidance throughout this process. Her comments helped me to further refine my bachelor thesis each time and gave me renewed motivation to continue working during difficult moments. I would also like to thank my co-supervisor, Mr. B. Oloeriu, for his guidance and support during this process.

In addition, I would like to thank the gymnasts and coaches of OTV Nazareth for their cooperation. Thanks to the provided video footage, the proof-of-concept could be tested on real executions of the kip on bars. In particular, I would like to thank the coach who helped provide the video material and who gave feedback on parts of my bachelor thesis based on her gymnastics experience.

I would also like to thank Tom Drijvers, my colleague during the internship, for proofreading parts of my bachelor thesis and for the additional feedback. Finally, I would like to thank everyone who supported, helped or motivated me during this period.

For me, this bachelor thesis forms a first exploration of how AI and video analysis can contribute to sport-technical support. I hope that this work can form a useful basis for further research within this domain.
```

